\subsection{Requirements of the Solution}
\speclabel{AO3.1.4a}

To ensure \projectNameFormatted\ executes deterministically, and can be tested against real-world attacks, the hardware and software requirements must be specified and justified based on the computational requirements.

Rather than just listing the requirements, below they are split into the \textbf{host environment} (the machine being protected) and the \textbf{evaluation device} (required to simulate attacks safely in a controlled environment).

\subsubsection{Hardware}

Table~\ref{tab:hardware_requirements} lists the hardware configurations required for deployment and testing.

	{\small
		\begin{xltabular}{\textwidth}{@{} p{2.8cm} p{2.4cm} L L @{}}
			\caption{Hardware Requirements and Justifications} \label{tab:hardware_requirements} \\
			\toprule
			\textbf{Component} & \textbf{Minimum Specification} & \textbf{Role} & \textbf{Justification} \\
			\midrule
			\endfirsthead

			\toprule
			\textbf{Component} & \textbf{Minimum Specification} & \textbf{Role} & \textbf{Justification} \\
			\midrule
			\endhead

			\midrule
			\multicolumn{4}{r@{}}{\scriptsize\textit{Continued on next page\dots}} \\
			\endfoot

			\bottomrule
			\endlastfoot

			\textbf{CPU} &
			x86-64 architecture; $\ge 2.0\,\text{GHz}$ base clock. &
			Runs the OS hook, heuristic engine, and UI threads. &
			Continuous rolling calculations require adequate hardware to prevent stalls during rapid bursts. \\
			\addlinespace[0.8em]

			\textbf{RAM} &
			$\ge 50\,\text{MB}$ free memory (system total $\ge 4\,\text{GB}$). &
			Stores the in-memory whitelists, rolling buffers, and thread states. &
			The buffers must be in physical RAM rather than the pagefile on disk, so the execution is sub-millisecond without causing I/O delays. \\
			\addlinespace[0.8em]

			\textbf{Secondary Storage} &
			$\ge 100\,\text{MB}$ available disk space. &
			Stores the compiled binary, configuration files, and event logs. &
			While the compiled executable requires minimal storage, persistent logs require sufficient headroom to prevent write failures. \\
			\addlinespace[0.8em]

			\textbf{USB Ports} &
			At least two independent physical ports. &
			Allows the simultaneous connection of a legitimate keyboard and an untrusted test peripheral. &
			Essential for making sure the OS assigns discrete hardware handles and the program identifies devices correctly. \\
			\addlinespace[0.8em]

			\textbf{Legitimate Peripheral} &
			Standard keyboard. &
			Used to gather standard human typing data during development and testing. &
			Required to collect inter-keystroke interval ($\text{IKI}$) benchmarks and make sure that legitimate inputs are not flagged. \\
			\addlinespace[0.8em]

			\textbf{BadUSB Test Device} &
			Programmable microcontroller (e.g., RP2040). &
			Simulates an attack in a controlled test environment. &
			Enables programmatic injection of known inputs at exact timings without deploying hazardous real-world malware. \\
		\end{xltabular}
	}

\begin{note}
	Using a dedicated programmable microcontroller (such as an RP2040) was chosen as a safety precaution. It allows a BadUSB exploit to be replicated, without exposing the testing machine to any actual danger.
\end{note}

\subsubsection{Software}

The software dependencies and toolchain have been chosen to ensure reliable execution speed and minimise the attack surface of the software itself.

	{\small
		\begin{xltabular}{\textwidth}{@{} p{2.8cm} p{2.4cm} L L @{}}
			\caption{Software Requirements and Justifications} \label{tab:software_requirements} \\
			\toprule
			\textbf{Component} & \textbf{Specification} & \textbf{Role} & \textbf{Justification} \\
			\midrule
			\endfirsthead

			\toprule
			\textbf{Component} & \textbf{Specification} & \textbf{Role} & \textbf{Justification} \\
			\midrule
			\endhead

			\midrule
			\multicolumn{4}{r@{}}{\scriptsize\textit{Continued on next page\dots}} \\
			\endfoot

			\bottomrule
			\endlastfoot

			\textbf{Operating System} &
			Microsoft Windows 10 / 11 (64-bit). &
			Target runtime environment. &
			\projectNameFormatted\ uses the Win32 APIs which are standardised and stable across all modern 64-bit Windows releases. \\
			\addlinespace[0.8em]

			\textbf{Programming Language} &
			ISO C. &
			The language used for implementing the daemon and heuristic engine. &
			Provides direct memory control with no garbage collection overhead, ensuring hooks execute well within their timeout limits. \\
			\addlinespace[0.8em]

			\textbf{Compiler} &
			Clang / LLVM. &
			Builds native binaries. &
			Produces standalone \texttt{.exe} binaries without requiring the .NET runtime or redistributables. \\
			\addlinespace[0.8em]

			\textbf{Win32 Core APIs} &
			\texttt{user32.dll} and \texttt{kernel32.dll}. &
			Provide hook installation, thread synchronisation, and Raw Input capabilities. &
			Standard system libraries linked at compile time, guaranteeing no external dependencies. \\
			\addlinespace[0.8em]

			\textbf{Data Serialiser} &
			Lightweight JSON parser (e.g., \texttt{cJSON}). &
			Serialises and deserialises the JSON files. &
			Allows human-readable persistent storage without introducing the bloat and attack surface of a fully fledged SQL database. \\
			\addlinespace[0.8em]

			\textbf{Benchmarking \& Analysis Tools} &
			Win32 High-Resolution Timers, Task Manager, and GitHub Actions. &
			Measure the memory footprint, CPU utilisation, and hook latency. &
			Required to quantitatively evaluate success criteria during testing. \\
		\end{xltabular}
	}
